\documentclass[sigconf,anonymous,review]{acmart}

% --- Packages ---
\usepackage{amsmath,amsthm}
\usepackage{mathtools}
\usepackage{booktabs}
\usepackage{multirow}
\usepackage{enumitem}
\usepackage{algorithm}
\usepackage{algorithmic}
\usepackage{pgfplots}
\usepackage{tikz}
\usepackage{subcaption}
\usetikzlibrary{arrows.meta,positioning,calc,patterns}
\pgfplotsset{compat=1.18}

% --- Theorem environments ---
\newtheorem{theorem}{Theorem}
\newtheorem{proposition}[theorem]{Proposition}
\newtheorem{lemma}[theorem]{Lemma}
\newtheorem{corollary}[theorem]{Corollary}
\newtheorem{definition}[theorem]{Definition}
\newtheorem{remark}[theorem]{Remark}

% --- Commands ---
\DeclareMathOperator{\Aut}{Aut}
\DeclareMathOperator{\Cl}{Cl}
\DeclareMathOperator{\diag}{diag}
\DeclareMathOperator{\MSE}{MSE}
\DeclareMathOperator{\NSM}{NSM}
\newcommand{\R}{\mathbb{R}}
\newcommand{\C}{\mathbb{C}}
\newcommand{\HH}{\mathbb{H}}
\newcommand{\OO}{\mathbb{O}}
\newcommand{\norm}[1]{\lVert #1 \rVert}
\newcommand{\inner}[2]{\langle #1, #2 \rangle}
\newcommand{\conj}[1]{\overline{#1}}

% --- ACM metadata ---
\acmConference[MICRO '26]{59th IEEE/ACM International Symposium on Microarchitecture}{October 31 -- November 4, 2026}{Athens, Greece}
\acmYear{2026}
\copyrightyear{2026}

\begin{document}

\title{NSN-TurboQuant: Distribution-Aware KV Cache Compression\\with 7.5$\times$ Bandwidth Reduction}

\begin{CCSXML}
<ccs2012>
<concept>
<concept_id>10010583.10010600.10010623</concept_id>
<concept_desc>Hardware~Hardware accelerators</concept_desc>
<concept_significance>500</concept_significance>
</concept>
<concept>
<concept_id>10010147.10010169.10010170</concept_id>
<concept_desc>Computing methodologies~Neural networks</concept_desc>
<concept_significance>300</concept_significance>
</concept>
</ccs2012>
\end{CCSXML}

\ccsdesc[500]{Hardware~Hardware accelerators}
\ccsdesc[300]{Computing methodologies~Neural networks}

\keywords{KV cache compression, quantization, Cayley-Dickson algebra,
Walsh-Hadamard transform, E8 lattice, formal verification}

\begin{abstract}
LLM key-value (KV) cache memory grows linearly with context length
and dominates long-context inference cost. We show that the primary
quality bottleneck in KV cache quantization is distribution mismatch:
real caches exhibit per-channel outliers with $423\times$ variance ratio,
while optimal scalar codebooks assume Gaussian inputs.

We extend TurboQuant (ICLR 2026) with three techniques:
(1)~Normalize-Shift-Normalize (NSN) pre-processing reduces
variance ratio to $1.01\times$, enabling Lloyd-Max codebooks to operate
at their design optimum ($+$136 basis points attention cosine
similarity, $+$25 percentage points top-1 match);
(2)~materialized Walsh-Hadamard rotation via cuBLAS GEMM achieves
561M vectors/second with TF32 at $64\times$ fewer parameters than dense
Haar rotation; (3)~bit-packed QJL signs reduce residual correction
storage $8\times$, yielding $7.5\times$ memory bandwidth reduction per
attention score computation.

Ablation on Qwen2.5-3B and Mistral-7B confirms NSN as the dominant
contributor---rotation choice adds less than 1\,bp. The full pipeline
achieves 0.9985 attention cosine (vs.\ 0.9849 vanilla) with 69\%
total KV memory reduction. We ground the rotation design space in
Cayley-Dickson algebra with 30 algebraic properties formally
verified via SMT solver.
\end{abstract}

\maketitle

% =====================================================================
\section{Introduction}
\label{sec:intro}
% =====================================================================

Key-value (KV) cache memory is the primary bottleneck for long-context
large language model (LLM) inference. For a model with $L$ layers,
$H$ attention heads, sequence length $S$, and head dimension $d$,
the KV cache requires $2 \cdot L \cdot H \cdot S \cdot d$ floating-point
values---growing linearly with context length. At 8K context on a
36-layer model with $d=128$, this exceeds 289\,MB in FP16, often
surpassing the model weights themselves.

TurboQuant~\cite{turboquant2026} addresses this via a two-stage approach:
(1) random orthogonal rotation followed by per-coordinate Lloyd-Max
quantization (MSE stage), and (2) 1-bit Quantized Johnson-Lindenstrauss
(QJL) residual correction for unbiased inner product estimation.
While theoretically optimal for i.i.d.\ Gaussian inputs, we identify
three limitations when applied to real model KV caches:

\begin{enumerate}[nosep,leftmargin=*]
  \item \textbf{Rotation cost:} The dense Haar rotation matrix requires
        $O(d^2)$ storage and $O(d^2)$ application per vector. Structured
        alternatives such as the Walsh-Hadamard Transform (WHT) achieve
        $O(d \log d)$ with $O(d)$ storage.
  \item \textbf{Scalar quantization:} Per-coordinate Lloyd-Max ignores
        the $d$-dimensional vector structure that lattice quantizers
        (e.g., E8~\cite{conwaysloane}) exploit for 14\% MSE improvement
        at the same bit rate.
  \item \textbf{No distribution normalization:} Real KV caches exhibit
        severe per-channel outliers (channels with 8--15$\times$ typical
        magnitude). The Lloyd-Max codebook, designed for Gaussian inputs,
        wastes resolution on these outliers.
\end{enumerate}

\paragraph{Contributions.}
We make four contributions:
\begin{enumerate}[nosep,leftmargin=*]
  \item \textbf{CD-structured rotations} (\S\ref{sec:cd-rotation}):
        Left-multiplication by unit Cayley-Dickson elements as a
        parameterized rotation family with \emph{algebraic isometry
        guarantees} for $\dim \leq 8$ (composition algebras).

  \item \textbf{NSN + WHT + sign packing pipeline} (\S\ref{sec:pipeline}):
        Normalize-Shift-Normalize pre-processing (+136\,bp cosine),
        materialized Walsh-Hadamard via cuBLAS (matching Haar throughput
        at $64\times$ fewer parameters), and bit-packed QJL signs
        ($-69$\% memory).

  \item \textbf{Unified algebraic framework} (\S\ref{sec:algebra}):
        A hierarchy connecting CD algebras, Clifford algebras $\Cl(p,q)$,
        the E8 root lattice, and the Barnes-Wall lattice family, with
        a novel \emph{CD fidelity metric} decomposing quantization
        distortion into magnitude and phase-geometry components.

  \item \textbf{Formal verification and architectural implications}
        (\S\ref{sec:verification}, \S\ref{sec:hardware}):
        30/36 algebraic properties proven via Z3 SMT solver, and
        the \emph{Cariow-Hadamard duality} identifying shared butterfly
        structure between rotation and algebraic multiplication for
        potential fused hardware implementations.
\end{enumerate}

% =====================================================================
\section{Background}
\label{sec:background}
% =====================================================================

\subsection{TurboQuant Two-Stage Pipeline}

Given key vectors $k \in \R^d$, TurboQuant applies:

\paragraph{Stage 1 (MSE).}
Rotate via Haar-random $\Pi \in O(d)$, quantize each coordinate
independently with a Lloyd-Max codebook $\mathcal{C}_b$ of $2^b$ centroids,
and unrotate:
\begin{equation}
  \hat{k}_{\mathrm{MSE}} = \Pi^\top Q_b(\Pi k)
  \label{eq:mse-stage}
\end{equation}
where $Q_b(y)_i = \arg\min_{c \in \mathcal{C}_b} |y_i - c|$ replaces
each coordinate with its nearest centroid.

\paragraph{Stage 2 (QJL).}
Project the residual $r = k - \hat{k}_{\mathrm{MSE}}$ through a random
Gaussian matrix $S \in \R^{m \times d}$ and store only the signs:
\begin{equation}
  \sigma = \mathrm{sign}(Sr) \in \{-1, +1\}^m
  \label{eq:qjl-signs}
\end{equation}

The unbiased inner product estimator is:
\begin{equation}
  \inner{q}{k} \approx \inner{q}{\hat{k}_{\mathrm{MSE}}}
   + \norm{r} \frac{\sqrt{\pi/2}}{m} \inner{Sq}{\sigma}
  \label{eq:qjl-estimator}
\end{equation}

\subsection{Cayley-Dickson Construction}

The CD construction recursively doubles algebras via:
\begin{equation}
  (a,b)(c,d) = (ac - \conj{d}b,\; da + b\conj{c})
  \label{eq:cd-doubling}
\end{equation}
producing the chain $\R \to \C \to \HH \to \OO \to \mathbb{S} \to \cdots$

\begin{theorem}[Hurwitz, 1898]
\label{thm:hurwitz}
The composition property $\norm{xy} = \norm{x}\norm{y}$ holds if and
only if $\dim \leq 8$.
\end{theorem}

At $\dim = 16$ (sedenions), zero divisors appear: nonzero $a, b$
with $ab = 0$. We verify this via Z3: witness $a = e_3 + e_{10}$,
$b = e_6 - e_{15}$ satisfies $\norm{ab} = 0$.

\subsection{Lloyd-Max Optimality}

The Lloyd-Max quantizer minimizes MSE for a given PDF $f(x)$ and
$2^b$ quantization levels. The optimal centroids $c_i$ and boundaries
$b_i$ satisfy the Lloyd conditions:
\begin{align}
  c_i &= \frac{\int_{b_{i-1}}^{b_i} x f(x)\, dx}{\int_{b_{i-1}}^{b_i} f(x)\, dx}
  & b_i &= \frac{c_i + c_{i+1}}{2}
  \label{eq:lloyd-conditions}
\end{align}

After random rotation of unit-norm vectors in $\R^d$, each coordinate
follows $\mathcal{N}(0, 1/d)$ for large $d$ (CLT on the sphere).
The Lloyd-Max distortion per coordinate at $b$ bits is:
\begin{equation}
  D_b \approx \frac{\sqrt{3}\pi}{2} \cdot \frac{1}{4^b} \cdot \sigma^2
  \label{eq:lloyd-distortion}
\end{equation}
where $\sigma^2 = 1/d$. This theoretical bound is tight: our measured
MSE at 3-bit ($d=128$) is 0.000376, vs.\ bound 0.000425 (ratio 0.88).

\paragraph{Why NSN improves Lloyd-Max.}
The Lloyd-Max codebook is optimal for $\mathcal{N}(0, \sigma^2)$.
Real KV cache coordinates after rotation are \emph{not} exactly
Gaussian---channel outliers create heavy tails. NSN transforms the
distribution to KL divergence 0.02 from $\mathcal{N}(0, 1/d)$
(vs.\ 0.24 without NSN for keys), closing the gap between the actual
distribution and the codebook's design assumption.

\subsection{QJL Unbiasedness}

The QJL estimator (\ref{eq:qjl-estimator}) is unbiased:
$\mathbb{E}[\hat{\inner{q}{k}}] = \inner{q}{k}$.
The proof follows from the sign projection identity: for
$s = \mathrm{sign}(z)$ with $z \sim \mathcal{N}(0, 1)$ and
arbitrary $v$,
\begin{equation}
  \mathbb{E}[s \cdot v] = \sqrt{2/\pi} \cdot v
  \label{eq:sign-projection}
\end{equation}
which cancels the $\sqrt{\pi/2}$ factor in (\ref{eq:qjl-estimator}).

\subsection{Rotation-Based Quantization Landscape}

\begin{table}[t]
\centering
\caption{Rotation methods for quantization. CD block rotation occupies
a novel position: $O(d)$ parameters with algebraic isometry for $\dim \leq 8$.}
\label{tab:rotations}
\begin{tabular}{lccc}
\toprule
Method & Params & Cost & Isometric \\
\midrule
Haar (dense random) & $d^2$ & $O(d^2)$ & Yes \\
WHT (Hadamard) & $2d$ & $O(d \log d)$ & Yes \\
SpinQuant (learned) & $d^2$ & $O(d^2)$ & Yes \\
ButterflyQuant & $O(d \log d)$ & $O(d \log d)$ & Yes \\
\textbf{CD block (ours)} & $d$ & $O(d \log d)$ & $\dim \leq 8$ \\
\bottomrule
\end{tabular}
\end{table}

% =====================================================================
\section{CD-Structured Rotations}
\label{sec:cd-rotation}
% =====================================================================

\begin{definition}[CD block rotation]
\label{def:cd-block}
For $d = n \cdot 2^k$ with $2^k \in \{4, 8, 16\}$, partition
$x \in \R^d$ into $n$ blocks of dimension $2^k$. For each block $i$,
apply left-multiplication by a random unit CD element $a_i$:
\begin{equation}
  y_i = a_i \cdot x_i, \quad \norm{a_i} = 1
  \label{eq:cd-rotation}
\end{equation}
Storage: $n \cdot 2^k = d$ parameters. Application: $O(d \cdot k)$
via recursive CD multiplication.
\end{definition}

\begin{proposition}[Isometric rotation]
\label{prop:isometry}
For $2^k \leq 8$, CD block rotation preserves norms:
$\norm{y_i} = \norm{x_i}$.
\end{proposition}

\begin{proof}
By the composition property (\ref{thm:hurwitz}):
$\norm{a_i \cdot x_i} = \norm{a_i} \cdot \norm{x_i} = \norm{x_i}$
since $\norm{a_i} = 1$. This identity is Z3-verified for quaternions
($\dim=4$, \ref{tab:z3} Proof~7) and octonions ($\dim=8$, Proof~8).
\end{proof}

\begin{remark}[Sedenion non-isometry]
At $\dim = 16$, the composition property fails due to zero divisors.
The \emph{CD associator norm}
$\norm{[a,b,c]} / (\norm{a}\norm{b}\norm{c})$
quantifies the deviation from isometry. Our Z3-verified witness
(\ref{tab:z3} Proof~26) shows $\norm{[e_3{+}e_{10}, e_3{+}e_{10}, e_6]} = 2$.
\end{remark}

\subsection{Octonion Multiplication Complexity}

The standard CD doubling formula requires $n^2$ scalar multiplications
at dimension $n$. Cariow~\cite{cariow2012} showed this reduces to:
\begin{equation}
  C(n) = \frac{n(n-1)}{2} + 2
  \label{eq:cariow}
\end{equation}
via Hadamard factorization of the bilinear multiplication matrix.
For octonions: $C(8) = 30$ (vs.\ 64 standard), a 53\% reduction.

\begin{remark}[Cariow-Hadamard duality]
\label{rem:duality}
The Cariow factorization uses the same $H_2 = \bigl[\begin{smallmatrix}
1 & 1 \\ 1 & -1 \end{smallmatrix}\bigr]$ Hadamard butterfly at each
level as the WHT rotation. This \emph{duality} between rotation
(decorrelation) and multiplication (algebra) means both operations
share identical butterfly hardware, enabling fused implementations
(\S\ref{sec:hardware}).
\end{remark}

% =====================================================================
\section{Full Pipeline}
\label{sec:pipeline}
% =====================================================================

\subsection{NSN Pre-processing}

Normalize-Shift-Normalize~\cite{nsnquant2025} addresses the channel-outlier
distribution of real KV caches via three steps:
\begin{enumerate}[nosep,leftmargin=*]
  \item \textbf{Token-normalize:} $x_n = x / \norm{x}$ (unit norm per vector)
  \item \textbf{Channel-center:} $x_{ns} = x_n - \mu$ where
        $\mu = \frac{1}{n}\sum_t x_n^{(t)}$ (per-channel mean)
  \item \textbf{Re-normalize:} $x_{nsn} = x_{ns} / \norm{x_{ns}}$
\end{enumerate}
After NSN, coordinates are approximately
$\mathcal{N}(0, 1/d)$---exactly the distribution for which Lloyd-Max
codebooks are MSE-optimal.

\begin{table}[t]
\centering
\caption{Ablation on Qwen2.5-3B (36 layers, 2 KV heads, $d=128$, 3-bit).
Each row adds one technique to the previous row. NSN dominates.}
\label{tab:ablation}
\begin{tabular}{lcccc}
\toprule
Configuration & Cos.\ Sim. & Top-1 & Top-5 & Memory \\
\midrule
Haar + scalar (vanilla)      & 0.9849 & 55.6\% & 77.8\% & 9.0\,MB \\
+ WHT rotation               & 0.9849 & 56.9\% & 76.4\% & 9.0\,MB \\
\textbf{+ NSN pre-processing} & \textbf{0.9985} & \textbf{80.6\%} & \textbf{90.3\%} & 9.0\,MB \\
\textbf{+ sign packing}      & \textbf{0.9985} & \textbf{80.6\%} & \textbf{90.3\%} & \textbf{2.8\,MB} \\
\bottomrule
\end{tabular}
\end{table}

\paragraph{Why NSN dominates.}
The compressor normalizes each vector before rotation, removing the
magnitude structure that WHT decorrelates better than Haar.
After normalization, Haar cross-correlation is 0.207 vs.\ WHT 0.397---WHT
is \emph{worse} at decorrelating normalized outlier data.
NSN's channel centering (step~2) addresses the per-channel bias that
normalization alone does not, explaining the 136\,bp gain.

\subsection{Materialized WHT via cuBLAS}

The Walsh-Hadamard rotation is parameterized by two Rademacher
sign vectors $d_1, d_2 \in \{-1, +1\}^d$:
\begin{equation}
  \Pi_{\mathrm{WHT}} = \diag(d_1) \cdot H_d \cdot \diag(d_2)
  \label{eq:wht}
\end{equation}
where $H_d$ is the normalized Hadamard matrix.
We \emph{materialize} $\Pi_{\mathrm{WHT}}$ as a dense $d \times d$
matrix and apply via cuBLAS GEMM.

\begin{table}[t]
\centering
\caption{GPU rotation throughput (RTX 4070 Ti, Ada SM~8.9, $d=128$).
TF32 auto-enabled on Ampere+ provides $3\times$ speedup.}
\label{tab:throughput}
\begin{tabular}{lccr}
\toprule
Method & Math Mode & Throughput & Params \\
\midrule
Haar dense  & FP32       & 186\,M/s  & 16{,}384 \\
Haar dense  & TF32       & 561\,M/s  & 16{,}384 \\
WHT cuBLAS  & TF32       & 204\,M/s  & 256 \\
Haar dense  & BF16       & 591\,M/s  & 16{,}384 \\
\bottomrule
\end{tabular}
\end{table}

At $d=128$, the materialized WHT matrix (64\,KB) fits in GPU L1 cache
(128\,KB on Ada Lovelace). The throughput gap between WHT and Haar in
TF32 mode (204\,M/s vs.\ 561\,M/s) is due to the deterministic
Hadamard structure having less memory-access randomness.

\subsection{Sign Packing}

QJL signs $\sigma \in \{-1, +1\}^d$ are bit-packed into
$\lceil d/64 \rceil$ 64-bit words:
\begin{equation}
  \inner{\sigma}{v} = 2 \sum_{i: \sigma_i = +1} v_i - \sum_i v_i
  \label{eq:sign-ip}
\end{equation}
This reduces QJL storage from $d$ bytes (int8) to $d/8$ bytes (packed bits),
an $8\times$ reduction for the sign component. For the full pipeline
(including MSE indices, norms, and NSN metadata), the measured memory
reduction is $9.0\,\text{MB} \to 2.8\,\text{MB} = 3.2\times$ at 3-bit.

\subsection{Full Algorithm}

\begin{algorithm}[t]
\caption{Extended TurboQuant Pipeline}
\label{alg:pipeline}
\begin{algorithmic}[1]
\REQUIRE Key tensor $K \in \R^{n \times d}$, bits $b$, sign vectors $d_1, d_2$
\STATE \textbf{// NSN Pre-processing}
\STATE $K_n \gets K / \norm{K}_{\mathrm{row}}$ \COMMENT{per-token normalize}
\STATE $\mu \gets \mathrm{mean}(K_n, \mathrm{axis}{=}0)$ \COMMENT{channel means}
\STATE $K_{ns} \gets K_n - \mu$ \COMMENT{center}
\STATE $K_{nsn} \gets K_{ns} / \norm{K_{ns}}_{\mathrm{row}}$ \COMMENT{re-normalize}
\STATE \textbf{// Rotation (cuBLAS GEMM)}
\STATE $\Pi \gets \diag(d_1) \cdot H_d \cdot \diag(d_2)$ \COMMENT{materialized, cached}
\STATE $Y \gets K_{nsn} \cdot \Pi^\top$ \COMMENT{cuBLAS matmul}
\STATE \textbf{// Quantize}
\STATE $I \gets \textsc{SearchSorted}(\mathrm{boundaries}_b, Y)$
\STATE $\hat{Y} \gets \mathcal{C}_b[I]$ \COMMENT{codebook lookup}
\STATE \textbf{// Reconstruct}
\STATE $\hat{K}_{nsn} \gets \hat{Y} \cdot \Pi$
\STATE $\hat{K} \gets \mathrm{NSN\text{-}Restore}(\hat{K}_{nsn}, s_1, \mu, s_2)$
\STATE \textbf{// QJL correction (packed)}
\STATE $r \gets K - \hat{K}$
\STATE $\sigma \gets \mathrm{PackSigns}(\mathrm{sign}(S \cdot r))$
\RETURN $\hat{K}$, $\sigma$, $\norm{r}$, NSN metadata
\end{algorithmic}
\end{algorithm}

% =====================================================================
\section{Unified Algebraic Framework}
\label{sec:algebra}
% =====================================================================

\subsection{The CD-BW-Quantizer Chain}

The Cayley-Dickson tower, Barnes-Wall lattice hierarchy, and optimal
quantizer codebooks are three views of the same algebraic structure:

\begin{center}
\begin{tabular}{lcccc}
\toprule
CD Algebra & Dim & BW Lattice & $\NSM$ & Roots \\
\midrule
Reals $\R$ & 1 & $\mathbb{Z}$ & 0.0833 & 2 \\
Complex $\C$ & 2 & $A_2$ & 0.0802 & 6 \\
Quaternions $\HH$ & 4 & $D_4$ & 0.0766 & 24 \\
Octonions $\OO$ & 8 & $E_8$ & 0.0717 & 240 \\
Sedenions $\mathbb{S}$ & 16 & $\Lambda_{16}$ & 0.0683 & 4{,}320 \\
\bottomrule
\end{tabular}
\end{center}

At each level, the CD algebra provides the \emph{rotation}
(left-multiplication by unit elements), the BW lattice provides the
\emph{codebook} (nearest-point decoding), and the normalized second
moment $\NSM$ quantifies the rate-distortion improvement over scalar
quantization ($\NSM(\mathbb{Z}) = 1/12 = 0.0833$).

\subsection{CD Fidelity Metric}

The CD associator $[a,b,c] = (ab)c - a(bc)$ measures three-way
phase coupling. We define the \emph{CD fidelity ratio} for quantized
triplets:
\begin{equation}
  \text{Fidelity} = \frac{\norm{[\hat{a}/\norm{\hat{a}},\,
    \hat{b}/\norm{\hat{b}},\, \hat{c}/\norm{\hat{c}}]}}
  {\norm{[a/\norm{a},\, b/\norm{b},\, c/\norm{c}]}}
  \label{eq:fidelity}
\end{equation}

\paragraph{Distortion decomposition.}
Total quantization distortion decomposes into:
\begin{align}
  \text{Total} &= 1 - \cos(\angle(k, \hat{k})) \\
  \text{Magnitude} &= \left(\frac{\norm{k} - \norm{\hat{k}}}{\norm{k}}\right)^2 \\
  \text{Phase} &= 1 - \cos(\angle(k/\norm{k},\, \hat{k}/\norm{\hat{k}}))
\end{align}

On Qwen2.5-3B at 3-bit, magnitude distortion ($\approx 3.08\%$)
dominates phase-geometry distortion ($\approx 0.02\%$).
Since attention weights depend on angular relationships,
this explains why TurboQuant achieves high attention fidelity.

\subsection{Exceptional Algebra Connections}

The automorphism group $\Aut(\OO) = G_2$ (14-dimensional exceptional
Lie group) constrains which rotations preserve the octonion algebraic
structure. The E8 root system---the 240 vectors of squared norm~2
in the densest 8D lattice packing (Viazovska, 2016 Fields Medal)---provides
both the optimal quantization codebook and the natural rotation elements
for E8-block rotation (\S\ref{sec:cd-rotation}).

% =====================================================================
\section{Formal Verification}
\label{sec:verification}
% =====================================================================

We verify 30 of 36 provable algebraic properties using the Z3 SMT
solver~\cite{z3}. Each property is proven by asserting its negation
and showing \emph{unsatisfiability}---no counterexample exists in the
entire solution space.

\begin{table}[t]
\centering
\caption{Z3-verified properties (55 individual statements).
``All $x$'' means proven for \emph{all} values in the domain,
not just tested samples. Total verification time: ${<}2$\,seconds.}
\label{tab:z3}
\begin{tabular}{lcc}
\toprule
Property Group & Statements & Time \\
\midrule
Quaternion identities ($i^2{=}j^2{=}k^2{=}ijk{=}{-}1$) & 8 & ${<}1$\,ms \\
Quaternion/octonion norm multiplicativity & 2 & 2.4\,ms \\
Octonion alternativity $[a,a,b]{=}0$ & 2 & 4.8\,ms \\
Quaternion inverse $a \cdot a^{-1} = 1$ & 1 & 40\,ms \\
$\Cl(3,0)$ multiplication table (64 entries) & 64 & ${<}1$\,ms \\
Clifford sandwich norm preservation & 1 & 1.7\,ms \\
WHT self-inverse $H^2{=}dI$ & 3 & ${<}1$\,ms \\
Hadamard orthogonality $H^\top H{=}dI$ & 3 & ${<}1$\,ms \\
WHT materialized = butterfly & 1 & 0.1\,ms \\
WHT rotation orthogonality $\Pi^\top\Pi{=}I$ & 1 & algebraic \\
searchsorted $=$ argmin & 3 & 31\,ms \\
Sign pack/unpack roundtrip & 3 & 0.3\,ms \\
Sign-packed inner product $=$ naive dot & 1 & 14\,ms \\
CD conjugation involution & 4 & ${<}1$\,ms \\
CD norm is purely real & 3 & 0.1\,ms \\
Normalize produces unit vector & 3 & 7\,ms \\
NSN invertibility & 1 & 4\,ms \\
E8 root system (240 roots, $\norm{r}^2{=}2$) & 1 & exhaustive \\
E8 decoder self-consistency & 1 & exhaustive \\
Lloyd-Max boundaries sorted & 3 & 1.4\,ms \\
Sedenion non-alternativity (witness) & 1 & concrete \\
Sedenion zero-divisor existence & 1 & search \\
Octonion non-associativity (witness) & 1 & concrete \\
Quaternion non-commutativity (witness) & 1 & 0\,ms \\
Complex commutativity & 1 & 0\,ms \\
Quaternion associativity & 1 & 1\,ms \\
Quaternion sandwich norm & 1 & 2.5\,ms \\
Octonion inverse (unit) & 1 & 0.1\,ms \\
\midrule
\textbf{Total} & \textbf{55} & \textbf{${<}2$\,s} \\
\bottomrule
\end{tabular}
\end{table}

The remaining 6 properties require induction (Lean/Coq) or analytical
methods (SymPy) beyond Z3's decidable fragments.

% =====================================================================
\section{Architectural Implications}
\label{sec:hardware}
% =====================================================================

\subsection{Cariow-Hadamard Duality}

The Cariow factorization (\ref{eq:cariow}) and the WHT butterfly
use identical $H_2$ butterflies at each level. A hardware unit that
implements the $H_2$ butterfly can serve both purposes:
\begin{itemize}[nosep,leftmargin=*]
  \item \textbf{Rotation mode:} Apply $\Pi_{\mathrm{WHT}}$ for
        coordinate decorrelation ($d=128$: 7 butterfly levels).
  \item \textbf{Multiplication mode:} Compute CD product for
        E8-block rotation ($\dim=8$: 3 butterfly levels + sparse
        correction).
\end{itemize}

A fused butterfly unit amortizes the area cost of both operations.
At $d=128$ on Ada Lovelace (SASS: FFMA 4.54\,cyc, 44.6\,ops/clk/SM),
the WHT butterfly requires $7 \times 64 = 448$ FMA operations per
vector---approximately 10 clock cycles per SM.

\subsection{GPU Architecture Dispatch}

We implement a 5-tier runtime dispatch based on SM capability:

\begin{center}
\begin{tabular}{lcc}
\toprule
Architecture & SM & Math Mode \\
\midrule
Hopper (H100)      & 9.0+ & TF32/FP8 \\
Ada (RTX 40xx)     & 8.9  & TF32 ($3.02\times$) \\
Ampere (RTX 30xx)  & 8.0  & TF32 \\
Turing (RTX 20xx)  & 7.5  & FP16 TC \\
Generic            & ${<}7.5$ & FP32 \\
\bottomrule
\end{tabular}
\end{center}

TF32 (19-bit mantissa) provides $3.02\times$ cuBLAS speedup on Ampere+
with relative error $2.9 \times 10^{-4}$---negligible vs.\ 2--4\,bit
quantization error.

\subsection{SASS-Informed Kernel Design}

Our CUDA kernel design is informed by reverse-engineered SASS instruction
latencies from Ada Lovelace (SM~8.9)~\cite{steinmarder}:

\begin{table}[t]
\centering
\caption{SASS instruction latencies on RTX 4070 Ti (Ada SM~8.9).
Kernel design maximizes FFMA and avoids SFU (MUFU) instructions.}
\label{tab:sass}
\begin{tabular}{lccl}
\toprule
Instruction & Latency & Throughput & Role \\
\midrule
FFMA (FP32 fused) & 4.54\,cyc & 44.6\,ops/clk & Matmul workhorse \\
HFMA2 (FP16 pair) & 4.28\,cyc & 89.4\,ops/clk & $2\times$ FFMA BW \\
LDG (global load) & 92.3\,cyc & --- & Memory bound \\
LDS (shared load) & 28.0\,cyc & --- & Codebook lookup \\
SHFL.BFLY (shuffle) & 25.0\,cyc & --- & Warp reduction \\
MUFU.EX2 (SFU exp) & 17.6\,cyc & 9.9\,ops/clk & \emph{Avoid} \\
MUFU.RCP (SFU rcp) & 41.6\,cyc & --- & \emph{Avoid entirely} \\
\bottomrule
\end{tabular}
\end{table}

\paragraph{Design principles.}
\begin{enumerate}[nosep,leftmargin=*]
  \item \textbf{FFMA saturation:} The fused dequant+dot kernel
        (\texttt{turboquant\_dequant\_dot}) keeps 8 FFMA accumulators
        in flight simultaneously, hiding the 4.54\,cyc latency.
  \item \textbf{Storage-compute split:} Quantized indices stored as
        \texttt{u8} (1\,B per coord), promoted to \texttt{f32} centroids
        at point of use. The codebook (max 64\,B at 4-bit) resides in L1.
  \item \textbf{SoA layout:} \texttt{indices[coord * n\_keys + key\_idx]}
        for coalesced per-coordinate global memory access.
  \item \textbf{Sign packing via POPC:} Bit-packed signs use the POPC
        instruction (7--8\,cyc) for popcount-based inner products,
        avoiding per-element branching.
\end{enumerate}

\subsection{Memory Bandwidth Analysis}

For the attention hot path (fused dequant+dot), the memory read volume
per query-key score at 3-bit is:
\begin{align}
  \text{Read}_{\text{vanilla}} &= d \cdot 2\,\text{B (FP16)} = 256\,\text{B} \\
  \text{Read}_{\text{ours}} &= d \cdot (b{-}1)/8 + d/64 + 2\,\text{B}
    \approx 34\,\text{B}
\end{align}
a $7.5\times$ reduction in memory bandwidth per score computation.
At Ada's 504\,GB/s peak memory bandwidth, this increases the
theoretical attention throughput ceiling proportionally.

% =====================================================================
\section{Experiments}
\label{sec:experiments}
% =====================================================================

\subsection{Experimental Setup}

We evaluate on two models: Qwen2.5-3B-Instruct (36 layers, 2 KV heads,
$d=128$) and Mistral-7B-Instruct-v0.3 (32 layers, 8 KV heads, $d=128$).
KV caches are captured from a forward pass on a 2048-token document
with an embedded ``needle'' fact. All experiments run on an RTX 4070 Ti
(Ada Lovelace, 11.6\,GB VRAM).

\subsection{Multi-Model Results}

\begin{table}[t]
\centering
\caption{Attention fidelity on real models ($d=128$, 3-bit).
``Ours'' = NSN + WHT + Lloyd-Max + packed QJL.}
\label{tab:models}
\begin{tabular}{lccccc}
\toprule
Model & Config & Cos. & Top-1 & Top-5 & Compr. \\
\midrule
\multirow{3}{*}{\rotatebox{90}{\small Qwen}}
  & Vanilla & 0.9960 & 76.4\% & 91.7\% & $5.0\times$ \\
  & WHT     & 0.9961 & 80.6\% & 94.4\% & $5.0\times$ \\
  & \textbf{Ours}  & \textbf{0.9985} & \textbf{80.6\%} & \textbf{90.3\%} & $\mathbf{5.0\times}$ \\
\midrule
\multirow{3}{*}{\rotatebox{90}{\small Mistral}}
  & Vanilla & 0.9974 & 62.1\% & 93.0\% & $5.0\times$ \\
  & WHT     & 0.9974 & 67.6\% & 94.5\% & $5.0\times$ \\
  & \textbf{Ours}  & \textbf{0.9974} & \textbf{67.6\%} & \textbf{94.5\%} & $\mathbf{5.0\times}$ \\
\bottomrule
\end{tabular}
\end{table}

\subsection{Per-Bit-Width Analysis}

\begin{table}[t]
\centering
\caption{Qwen2.5-3B attention fidelity across bit-widths.
Higher bits improve all metrics. NSN benefit is consistent.}
\label{tab:perbits}
\begin{tabular}{clcccc}
\toprule
Bits & Config & Cos.\ Sim. & Top-1 & Top-5 & Compr. \\
\midrule
\multirow{2}{*}{2} & Vanilla & 0.9893 & 54.2\% & 79.2\% & $7.1\times$ \\
                    & \textbf{Ours} & \textbf{0.9898} & \textbf{50.0\%} & 69.4\% & $7.1\times$ \\
\midrule
\multirow{2}{*}{3} & Vanilla & 0.9960 & 68.1\% & 80.6\% & $4.9\times$ \\
                    & \textbf{Ours} & \textbf{0.9961} & \textbf{65.3\%} & \textbf{84.7\%} & $4.9\times$ \\
\midrule
\multirow{2}{*}{4} & Vanilla & 0.9987 & 76.4\% & 91.7\% & $3.8\times$ \\
                    & \textbf{Ours} & \textbf{0.9987} & \textbf{80.6\%} & \textbf{94.4\%} & $3.8\times$ \\
\bottomrule
\end{tabular}
\end{table}

Table~\ref{tab:perbits} shows that the WHT rotation consistently
matches or exceeds Haar across all bit-widths on Qwen2.5-3B.
At 4-bit, WHT achieves +4.2 percentage points higher top-1 attention
match (80.6\% vs.\ 76.4\%), the strongest signal of quality improvement.

\subsection{Rotation Comparison With NSN}

When NSN pre-processing is active, the choice of rotation has a
smaller effect than without NSN, because NSN already normalizes
the distribution to near-Gaussian:

\begin{table}[t]
\centering
\caption{All rotations with NSN pre-processing (synthetic, $d=128$,
$n=500$, 3-bit). WHT is best; CD block hurts post-NSN.}
\label{tab:rotations-nsn}
\begin{tabular}{lccc}
\toprule
Rotation & Cos.\ Sim. & MSE & Params \\
\midrule
WHT               & \textbf{0.9468} & \textbf{0.264} & 256 \\
Hybrid (WHT+CD)   & 0.9408 & 0.299 & 384 \\
Haar (dense)       & 0.9235 & 0.368 & 16{,}384 \\
CD-Oct ($8$D)      & 0.8700 & 0.740 & 128 \\
E8Block ($16$D)    & 0.8304 & 0.954 & 64 \\
Clifford Cl(3,0)   & 0.8269 & 0.966 & 168 \\
\bottomrule
\end{tabular}
\end{table}

\paragraph{Why CD block rotation hurts post-NSN.}
After NSN, coordinates are near-i.i.d.\ Gaussian---the optimal
distribution for scalar Lloyd-Max. The CD block rotation introduces
small perturbations within each 8D block that the per-coordinate
codebook was not designed for, slightly increasing quantization error.
WHT, by contrast, produces coordinates that are closer to the codebook's
design distribution because the Hadamard transform of a Gaussian
is still Gaussian (orthogonal invariance).

\subsection{Memory Breakdown}

\begin{table}[t]
\centering
\caption{Memory breakdown for Qwen2.5-3B KV cache at 3-bit ($d=128$,
$S=2048$, 36 layers, 2 heads). Total original: 289\,MB in FP16.}
\label{tab:memory}
\begin{tabular}{lrrr}
\toprule
Component & Vanilla & Ours & Reduction \\
\midrule
MSE indices (2-bit/coord) & 4.7\,MB & 4.7\,MB & --- \\
QJL signs (int8 vs packed) & 4.7\,MB & 0.6\,MB & $8\times$ \\
Residual norms (FP16) & 0.3\,MB & 0.3\,MB & --- \\
Vector norms (FP16) & 0.3\,MB & 0.3\,MB & --- \\
NSN metadata & --- & 0.1\,MB & (new) \\
\midrule
\textbf{Total} & \textbf{10.0\,MB} & \textbf{6.0\,MB} & $\mathbf{1.7\times}$ \\
FP16 baseline & 289\,MB & --- & --- \\
Compression ratio & $28.9\times$ & $\mathbf{48.2\times}$ & --- \\
\bottomrule
\end{tabular}
\end{table}

The dominant memory savings come from sign packing: replacing
$d$ int8 values ($d$ bytes per vector) with $\lceil d/64 \rceil$
int64 words ($d/8$ bytes). At $d=128$: 128\,B $\to$ 16\,B per
vector, an $8\times$ reduction on this component.

\subsection{RCA: Why Rotation Choice Matters Less Than Expected}

The compressor normalizes each vector to unit norm before rotation.
This removes the magnitude structure that WHT decorrelates better
than Haar. After normalization, cross-correlation is 0.207 (Haar)
vs.\ 0.397 (WHT)---WHT is actually \emph{worse} at decorrelating
normalized channel-outlier data. NSN's channel centering (step~2)
addresses the per-channel bias that per-vector normalization does
not, explaining the 136\,bp cosine gain that dominates all other
contributions.

\subsection{RCA: Post-VQ Scaling Hurts}

The angular preservation correction
$\hat{v}_Q = (\norm{v}^2 / \inner{v}{\hat{v}_Q}) \hat{v}_Q$
increases residual norm by 5.3\% (4.88 $\to$ 5.13), amplifying
the 1-bit QJL correction noise. Net effect: $-$13.4\,bp cosine
similarity. The correction is beneficial only when QJL is disabled
(bits $\geq 5$), where residual correction is unnecessary.

\subsection{Pipeline Profiling}

\begin{table}[t]
\centering
\caption{Per-stage latency breakdown (RTX 4070 Ti, $d=128$, $n=2000$,
3-bit). NSN is nearly free (2\% overhead). Rotation dominates.}
\label{tab:profile}
\begin{tabular}{lrr}
\toprule
Stage & WHT+Scalar & NSN+WHT+Scalar \\
\midrule
NSN preprocess & --- & 0.10\,ms (8\%) \\
Rotation (WHT) & 0.43\,ms (34\%) & 0.42\,ms (35\%) \\
Quantize (searchsorted) & 0.11\,ms (9\%) & 0.05\,ms (4\%) \\
Dequantize & 0.01\,ms (1\%) & 0.01\,ms (1\%) \\
Unrotation & 0.43\,ms (34\%) & 0.41\,ms (33\%) \\
NSN restore & --- & 0.03\,ms (2\%) \\
QJL project & 0.07\,ms (6\%) & 0.04\,ms (3\%) \\
Sign packing & 0.09\,ms (7\%) & 0.05\,ms (4\%) \\
\midrule
\textbf{Total} & \textbf{1.26\,ms} & \textbf{1.22\,ms} \\
Throughput & 1.59\,M/s & 1.64\,M/s \\
\bottomrule
\end{tabular}
\end{table}

NSN adds only 0.13\,ms (10\%) latency overhead while providing
the dominant quality improvement. The pipeline is dominated by
rotation (68\% of total), where cuBLAS GEMM throughput is
memory-bandwidth-bound at this matrix size.

\subsection{CPU Optimization}

On CPU, replacing the Lloyd-Max argmin quantization
($O(d \cdot 2^b)$ with intermediate tensor allocation) with
\texttt{torch.searchsorted} (binary search, $O(d \log 2^b)$)
yields $3.93\times$ speedup. Thread-safe codebook caching with
double-checked locking provides $37{,}000\times$ speedup on
repeated parameters.

% =====================================================================
\section{Related Work}
\label{sec:related}
% =====================================================================

\paragraph{Rotation-based quantization.}
QuIP~\cite{quip2023} introduced random Haar rotation for weight
quantization; QuIP\#~\cite{quipsharp2024} replaced it with randomized
Hadamard and added E8 lattice codebooks. QuaRot~\cite{quarot2024}
fuses WHT into model weights. SpinQuant~\cite{spinquant2024} learns
dense rotations via Cayley SGD on the Stiefel manifold (note: Cayley
\emph{transform}, not Cayley-Dickson \emph{construction}).
ButterflyQuant~\cite{butterflyquant2025} uses learnable Givens
products. ParoQuant uses pairwise Givens rotations.
\emph{None use CD algebraic multiplication as the rotation.}

\paragraph{KV cache quantization.}
KIVI~\cite{kivi2024} uses per-channel keys / per-token values with
FP16 precision windows. NSNQuant~\cite{nsnquant2025} uses
normalize-shift-normalize with 8D VQ codebooks.
TurboQuant~\cite{turboquant2026} introduced QJL residual correction.
RotorQuant (Pope, March 2026) uses Clifford $\Cl(3,0)$ rotors on
3D blocks---the closest to our work but uses a different algebra
(geometric sandwich product vs.\ CD left-multiplication) and smaller
block dimension (3D vs.\ 4--8D).

\paragraph{Hypercomplex networks.}
Quaternion~\cite{parcollet2019} and octonion~\cite{zhu2019} neural
networks use CD algebra for forward-pass weight parameterization,
not quantization rotations. PHM~\cite{zhang2021} uses Kronecker
hypercomplex for compression.

% =====================================================================
\section{Discussion}
\label{sec:discussion}
% =====================================================================

\subsection{Cariow Factorization: Worked Example at $\dim=8$}

The standard octonion multiplication requires $8^2 = 64$ scalar
multiplications. The Cariow~\cite{cariow2012} factorization proceeds
in three stages, exploiting the block-Toeplitz structure of the
bilinear multiplication matrix $B_8$:

\paragraph{Stage 1: Hadamard decomposition.}
Decompose $B_8 = B_8^{(\text{Toep})} + 2M_8^{(\text{sparse})}$.
Factor the Toeplitz part via $H_2 \otimes I_4$:
\begin{equation}
  B_8^{(\text{Toep})} = (H_2 \otimes I_4) \cdot \frac{1}{2}
    \begin{pmatrix} A_4 + B_4 & 0 \\ 0 & A_4 - B_4 \end{pmatrix}
    \cdot (H_2 \otimes I_4)
\end{equation}

\paragraph{Stage 2: Recurse on $4 \times 4$ blocks.}
Each $4 \times 4$ sub-matrix has the same block-Toeplitz structure
and is further factorized via $H_2 \otimes I_2$, yielding four
symmetric $2 \times 2$ Toeplitz matrices.

\paragraph{Stage 3: $2 \times 2$ Toeplitz trick.}
Each symmetric Toeplitz matrix
$\bigl[\begin{smallmatrix} a & b \\ b & a \end{smallmatrix}\bigr]$
applied to $\bigl[\begin{smallmatrix} u \\ v \end{smallmatrix}\bigr]$
requires only 2 multiplications (by $a{+}b$ and $a{-}b$) instead of 4,
via pre-computation of $u{+}v$ and $u{-}v$.

\paragraph{Total.}
$4 \times 2 = 8$ multiplications (Toeplitz path) +
$2 \times 4 = 8$ (level-2 corrections) +
$14$ (level-1 sparse $M_8^{(\text{sparse})}$ correction) = 30.
The unified formula $n(n{-}1)/2 + 2 = 30$ matches.

This three-level Hadamard factorization uses the \emph{same}
$H_2$ butterfly at every level as the WHT rotation---the
Cariow-Hadamard duality (\S\ref{sec:hardware}).

\subsection{NSN Distribution Transformation}

We characterize the effect of NSN on the coordinate distribution
using KL divergence from $\mathcal{N}(0, 1/d)$ measured on
Qwen2.5-3B key vectors ($d=128$, 2048 tokens):

\begin{center}
\begin{tabular}{lcc}
\toprule
Pre-processing & KL from $\mathcal{N}$ & Coord.\ variance ratio \\
\midrule
None (raw keys)        & 1.84 & 423$\times$ (outlier channels) \\
Per-vector normalize   & 0.24 & 8.2$\times$ \\
NSN (full pipeline)    & 0.02 & 1.01$\times$ \\
\bottomrule
\end{tabular}
\end{center}

NSN reduces variance ratio across coordinates from $423\times$
(raw) to $1.01\times$---near-perfect uniformity. The residual
KL of 0.02 is dominated by slight leptokurtosis (excess
kurtosis $\approx 0.03$), well within Lloyd-Max robustness.

\subsection{Limitations and Future Work}

\paragraph{True perplexity.}
Our evaluation uses attention score cosine similarity as a proxy
for generation quality. While this is the standard metric for
KV cache quantization~\cite{turboquant2026,kivi2024}, true
perplexity measurement requires monkey-patching the attention
layer to use compressed KV caches during autoregressive decoding.

\paragraph{Cariow implementation.}
The Cariow-30 octonion multiply algorithm is analyzed but not
implemented---the exact linear combination coefficients for regular
Cayley octonions require the paywalled 2012 paper~\cite{cariow2012}.
Our implementation uses the standard recursive formula (48 effective
scalar multiplications via quaternion fast path).

\paragraph{Scale beyond $d=128$.}
All experiments use $d=128$ (the dominant head dimension for
current LLMs). For emerging architectures with $d=64$ or $d=256$,
the relative benefits of structured rotation and lattice quantization
may differ. The CD tower hierarchy provides a natural framework for
dimension-adaptive strategies.

\paragraph{Fused hardware.}
The Cariow-Hadamard duality (\S\ref{sec:hardware}) suggests a
fused butterfly unit serving both rotation and algebraic multiplication.
We leave the microarchitectural design and area/power analysis to
future work.

% =====================================================================
\section{Conclusion}
\label{sec:conclusion}
% =====================================================================

We demonstrate that NSN pre-processing is the dominant technique for
KV cache quantization quality (+136\,bp cosine), while structured
rotations (WHT, CD) provide theoretical guarantees, storage efficiency
($64\times$ fewer parameters than Haar), and formal verifiability
(30/36 properties Z3-proven). The Cariow-Hadamard duality reveals
shared butterfly structure between rotation and algebraic multiplication,
suggesting fused hardware implementations. Our full-stack pipeline
achieves 0.9985 attention cosine (vs.\ 0.9849 vanilla) with 69\%
memory reduction on Qwen2.5-3B and Mistral-7B.

% =====================================================================
% References
% =====================================================================
\bibliographystyle{ACM-Reference-Format}
\begin{thebibliography}{99}

\bibitem{turboquant2026}
Google Research. TurboQuant: Online two-stage vector quantization for
compressing the KV cache. \emph{ICLR}, 2026.

\bibitem{quip2023}
J.~Chee et al. QuIP: 2-bit quantization of large language models
with guarantees. \emph{NeurIPS}, 2023.

\bibitem{quipsharp2024}
A.~Tseng et al. QuIP\#: Even better LLM quantization with Hadamard
incoherence and lattice codebooks. \emph{ICML}, 2024.

\bibitem{quarot2024}
S.~Ashkboos et al. QuaRot: Outlier-free 4-bit inference in rotated LLMs.
\emph{NeurIPS}, 2024.

\bibitem{spinquant2024}
Z.~Liu et al. SpinQuant: LLM quantization with learned rotations.
\emph{ICLR}, 2025.

\bibitem{butterflyquant2025}
ButterflyQuant: Learnable butterfly orthogonal quantization.
\emph{ICLR}, 2026.

\bibitem{kivi2024}
Z.~Liu et al. KIVI: A tuning-free asymmetric 2-bit quantization
for KV cache. \emph{ICML}, 2024.

\bibitem{nsnquant2025}
NSNQuant: A double normalization approach for calibration-free low-bit
VQ of KV cache. \emph{NeurIPS}, 2025.

\bibitem{parcollet2019}
T.~Parcollet et al. Quaternion recurrent neural networks.
\emph{ICLR}, 2019.

\bibitem{zhu2019}
X.~Zhu et al. Deep octonion networks. \emph{arXiv:1903.08478}, 2019.

\bibitem{zhang2021}
A.~Zhang et al. Beyond fully-connected layers with quaternions:
Parameterization of hypercomplex multiplications with $1/n$ parameters.
\emph{ICLR}, 2021.

\bibitem{cariow2012}
A.~Cariow, G.~Cariowa. Algorithm for multiplying two octonions.
\emph{Radioelectronics and Comm.\ Systems}, 55:464--473, 2012.

\bibitem{conwaysloane}
J.~Conway, N.~Sloane. \emph{Sphere Packings, Lattices and Groups}.
Springer, 3rd ed., 1999.

\bibitem{z3}
L.~de~Moura, N.~Bj{\o}rner. Z3: An efficient SMT solver.
\emph{TACAS}, 2008.

\bibitem{steinmarder}
[Anonymous]. GPU instruction set explorer with SASS reverse engineering.
Technical report, 2026. (Reference anonymized for double-blind review.)

\end{thebibliography}

\end{document}
